\chapter{Uvod}

Strojno učenje je v zadnjih letih postalo ključna tehnologija v industriji in raziskovalnem svetu.
Modeli, ki so bili nekoč zgolj akademski eksperimenti, danes poganjajo produkcijske sisteme, ki jih
dnevno uporablja na milijone uporabnikov. S tem prehodom v produkcijo je postala kritična potreba
po sistematičnem pristopu k razvoju, uvajanju in vzdrževanju modelov strojnega učenja.

MLOps (\textit{Machine Learning Operations}) je disciplina, ki združuje prakse strojnega učenja in
operativnega inženirstva programske opreme z namenom poenostavitve in avtomatizacije celotnega
življenjskega cikla modelov~\cite{Kreuzberger2023}. Ker gre za razmeroma mlado področje, se nabor
orodij in pristopov hitro razvija, kar otežuje pregled in primerjavo obstoječih rešitev.

V tem diplomskem delu bomo sistematično pregledali in primerjali tehnologije, ki se danes
najpogosteje uporabljajo pri uvajanju praks MLOps. Opisali bomo ključne koncepte in faze
življenjskega cikla modela, pregledali reprezentativna orodja za vsako fazo ter jih primerjali
glede na funkcionalnost, zrelost in primernost za različne organizacijske kontekste.
